\section{Unique Capabilities}

\label{sec:unique}
%% ============================================================================

Several capabilities of the Hanzo + Lux stack have no equivalent in any combination of alternatives:

\subsection{Post-Quantum Readiness}

The stack implements all three NIST post-quantum standards (FIPS~203, 204, 205) across the full path: key encapsulation in TLS handshakes (ZAP protocol), digital signatures in consensus (Quasar dual-certificate), key derivation in identity (IAM session keys), and key share distribution in MPC (Lux MPC). No competing stack provides post-quantum cryptography at any layer, let alone all layers.

\subsection{FHE-Native Confidential Computing}

The T-Chain provides on-chain FHE computation with threshold decryption governed by consensus. This enables use cases impossible on any other platform: encrypted auctions where bids are never revealed until settlement, confidential DeFi where trade amounts are hidden from validators, and private AI inference where model inputs remain encrypted throughout the pipeline.

\subsection{Zero-Allocation Wire Protocol}

ZAP's zero-allocation design eliminates garbage collection pressure in Go validators, enabling sustained 180K TPS without GC pauses. gRPC-based validators (Tendermint, CometBFT) experience 50--200ms GC pauses under equivalent load, causing finality jitter.

\subsection{Dual-Certificate Consensus}

Quasar's dual Ed25519 + ML-DSA-65 block signatures provide quantum resistance today without waiting for the ecosystem to migrate. Clients choose which signature to verify based on their security requirements. No other consensus protocol offers this hybrid approach.

\subsection{Unified AI + Blockchain Identity}

A single Hanzo IAM OIDC token authorizes: API requests through the gateway, database queries in Base, secret access in KMS, model inference in Engine, agent execution in Agent, and transaction signing on the blockchain. No other platform provides a single identity that spans both Web2 and Web3 operations.

%% ============================================================================
